\chapter{Conclusion}\label{chp:LABEL_CHP_5}

In this work, we demonstrated the versatility and expressiveness of the ray marching technique. Initially, we compared ray tracing and ray marching, highlighting their similarities and differences. Signed Distance Functions were also presented as a flexible mathematical representation of 
geometry, allowing us to easily apply techniques for the construction of procedural scenes. After building a strong foundation upon these concepts, we explored how this technique can be employed to render volumetric clouds in real time. We were effectively able to produce visually convincing results, while adopting an iterative approach that presented code snippets and algorithms supported by rendered examples, ultimately providing a comprehensive overview of ray marching and its capabilities.

\section{Future Works}

% Beer's Law & Powder Effect https://progmdong.github.io/2019-03-04/Volumetric_Rendering/
% Mencionar outros metodos de obter formas complexas como:
% GRIDS: https://renderdiagrams.org/2017/12/28/signed-distance-fields/
% https://steamcdn-a.akamaihd.net/apps/valve/2007/SIGGRAPH2007_AlphaTestedMagnification.pdf
% NEURAL NETWORKS: 